\paragraph{Reference scores and cohort curation.}
The cohort was re-curated from the institutional ultrasound database described by \citet{xu2026sfibai}. We retained its expert-assigned ultrasound appearance scores, ranging from 0.0 to 3.5 in increments of 0.1. The published annotation protocol refines four clinical grades into this score scale and describes standardized instruction, centralized review, and expert review of grading disagreements. The present study cleaned and organized the existing images and annotations without regrading them. Its cohort and splits differ from those of the earlier study; the SFibAI baseline was trained and evaluated on the present splits.

\paragraph{Acquisition views and lesion-location annotations.}
Six-class view labels were recovered from the plane identifiers recorded during image acquisition and organized during data cleaning. Lesion-location supervision likewise uses the existing clinician-drawn rectangular boxes described by \citet{xu2026sfibai}, after cleaning and organization. These boxes mark image regions requiring attention and carry the original grading annotations. View labels therefore describe the recorded acquisition plane, whereas boxes identify local regions within the image. Neither annotation type is required at inference.

\paragraph{Coverage and weak-target scope.}
The curated records contain view and lesion-location annotations for every image, although not every record supplies an eligible nonempty weak target. The box objective uses valid nonempty boxes with a positive target score, as specified in Appendix~\ref{sec:appendix-implementation}. Weak-localization evaluation includes 4,035 eligible test images; view evaluation includes all 4,107 test images. Local boxes do not delineate complete lesion contours, so responses beyond them require separate assessment of spatial relevance.
